\section{Two Stacks and a Preference}
\label{sec:ch16-two-stacks}
\index{WorkScheduler@\code{WorkScheduler}|(}%

The thirteenth middleware calls \code{\_queryExecutor.ExecuteAsync}, and that is
the line chapter~\ref{ch:the-life-of-a-request} stopped at. Below it is
\code{WorkScheduler}, one per operation context, and it is worth reading for two
reasons: it explains how much concurrency you are actually getting, and it is
where the last of chapter~\ref{ch:the-life-of-a-request}'s owed items turns out
to have a surprising answer.

\index{WorkQueue@\code{WorkQueue}!is a pair of stacks}%
The scheduler holds two queues, one for parallel work and one for serial work.
Neither is a queue, and \code{WorkQueue.cs} opens by saying so:

\begin{minted}[fontsize=\footnotesize,breaklines]{csharp}
internal sealed class WorkQueue
{
    private readonly Stack<IExecutionTask> _immediateStack = new();
    private readonly Stack<IExecutionTask> _deferredStack = new();
\end{minted}

\code{TryTake} pops the immediate stack and falls back to the deferred one, so
within a bucket the order is most-recently-pushed first. The only ordering
guarantee anywhere in here is that immediate work is taken before deferred work.
If you have ever wondered why sibling resolvers do not run in the order you
wrote them, that is the answer, and it is a stack rather than anything cleverer.

\index{WorkScheduler@\code{WorkScheduler}!parallel work is preferred}%
The drive loop prefers parallel work absolutely. A serial task is only dequeued
when the parallel queue is empty \emph{and} has nothing still running, one at a
time, and then the loop waits for that one task before going round again. The
comment says why: \enquote{Parallel work is always preferred, so we take a
single serial task and see if this results in more parallel work.} Which is how
the mutation rule from section~\ref{sec:ch16-pure-or-a-task} ends up meaning
what the specification wants. One root mutation field's entire subtree drains
out of the parallel queue before the next root mutation field starts, and the
serialisation is at the granularity of a root field rather than of a resolver.

\index{WorkScheduler@\code{WorkScheduler}!the buffer is not a limit}%
There is one number in \code{WorkScheduler.Execute.cs} that invites a wrong
conclusion, and I reached for it before reading what happens to it:

\begin{minted}[fontsize=\footnotesize,breaklines]{csharp}
private readonly IExecutionTask?[] _buffer = new IExecutionTask?[ProcessorCount * 2];
\end{minted}

That is not a concurrency limit. It bounds how many tasks are drained per
acquisition of the lock, to cut contention, and every non-serial task in the
batch is then started and not awaited: \enquote{if work is NOT serial we will
just enqueue it and not wait for it to finish.} HotChocolate does not cap how
many of your resolvers run at once. The thread pool does, and if you want a
ceiling, the concurrency gate chapter~\ref{ch:the-life-of-a-request} watched
remove itself is the setting that provides one.

\subsection*{The event that counts nothing}

\index{RunTask@\code{RunTask}!what it actually counts}%
Chapter~\ref{ch:the-life-of-a-request} recorded three things it had not
verified, and told itself to say nothing about task counts. The specific item
was \code{RunTask}, a diagnostic event sitting on the same listener as the
per-field event, gated by the same property, and never measured.

The call graph is three links long and every one of them is narrow. The event is
raised in exactly one place, in
\code{OperationContext.IExecutionTaskContext.cs}:

\begin{minted}[fontsize=\footnotesize,breaklines]{csharp}
IDisposable IExecutionTaskContext.Track(IExecutionTask task)
{
    AssertInitialized();
    return DiagnosticEvents.RunTask(task);
}
\end{minted}

\code{Track} is called from exactly one place, inside the base class
\code{ExecutionTask}. And exactly two types derive from \code{ExecutionTask}:
\code{DeferTask}, which orchestrates one \code{@defer} branch, and a private
\code{NoOpExecutionTask} that exists for the case where every root field was
skipped, with a comment saying so in as many words. The two types that actually
run your resolvers are \code{ResolverTask} and \code{BatchResolverTask}. Both
implement the task interface directly, and neither passes through that base
class at all.

So \code{RunTask} fires once per deferred branch, once in an edge case you will
never deliberately hit, and otherwise never. On a query with no \code{@defer} it
fires zero times, which is every query Mosaic has ever served, because
chapter~\ref{ch:real-time-in-a-federated-world} established that
\code{EnableDefer} defaults off and no subgraph turns it on.

\index{RunTask@\code{RunTask}!measured against a defer}%
That is a claim about an absence, so it wants a run where the number is not
zero. The sample turns \code{EnableDefer} on and asks the same executor for the
same data twice, once plainly and once with a deferred fragment:

\begin{minted}[fontsize=\footnotesize,breaklines]{text}
== run-task-counts-no-resolver
   plain      -> 4 resolver tasks, 0 execution tasks
   with defer -> execution tasks 3
   PASS
\end{minted}

Four resolver tasks and no execution tasks for the ordinary query. Same
listener, same flag, same schema: the event only has anything to count once
something defers.

\index{diagnostic events!one flag, two unrelated events}%
Two things about that are worth carrying away. The first is that the names point
the wrong way round. \code{ResolveFieldValue} sounds like the special case and
is the one that fires on every field that costs a task; \code{RunTask} sounds
like the general case and is the rare one. The second is that they share a
gate. Turning on the per-field property to get resolver timings also subscribes
you to defer orchestration, and the interface documents both as being controlled
by that one flag. Two unrelated event streams behind one switch is a thing to
know before you write a listener that assumes the switch means what it is named
after.

\index{RunTask@\code{RunTask}!the documented example that cannot happen}%
There is one more wrinkle, and it is the reason I chased this as far as I did.
The interface's own documentation offers, as its example of what a
\code{RunTask} might be, \enquote{DataLoader batch execution or other background
processing tasks}. DataLoader dispatch does not go through the execution-task
machinery at all. It is a separate component with diagnostic events of its own,
and it is the subject of the next section. The example in the documentation
describes no code path that raises this event.

\medskip
Which puts a piece back into an argument
chapter~\ref{ch:data-without-the-n-plus-1} left standing. That chapter read the
DataLoader dispatcher end to end and found that nothing in it observes resolver
readiness, which is true and which leaves an obvious question hanging: something
must, or the phrase would not be in the documentation at all. The scheduler is
the other side of that boundary.

\index{WorkScheduler@\code{WorkScheduler}|)}%
